\documentclass[10pt]{article}
\usepackage[margin=0.6in]{geometry}
\usepackage{amsmath,amssymb}
\usepackage{multicol}
\usepackage{booktabs}
\usepackage{array}
\usepackage{enumitem}
\usepackage[dvipsnames]{xcolor}

\pagestyle{empty}

% Tableau signs
\newcommand{\T}{\mathbb{T}}
\newcommand{\F}{\mathbb{F}}

% Vertical dots shorthand
\newcommand{\vmark}{\,\vert\,}

% Section formatting
\newcommand{\ruleheader}[1]{\medskip\noindent\textbf{\large #1}\par\smallskip\hrule\smallskip}
\newcommand{\rulesubheader}[1]{\smallskip\noindent\textbf{#1}\par\smallskip}

% Branching notation
\newcommand{\branches}[2]{%
  \begin{tabular}{@{}c@{\qquad}c@{}}
  #1 & #2
  \end{tabular}%
}

\setlength{\parindent}{0pt}
\setlength{\parskip}{2pt}

\begin{document}

\begin{center}
{\LARGE\bfseries Tableau Rules Cheat Sheet}\\[4pt]
{\large Phil 305 --- Propositional and Modal Logic}
\end{center}

\ruleheader{Part 1: Propositional Connectives}

Each connective has a rule for when the sentence is \textbf{true} ($\T$) and when it is \textbf{false} ($\F$). Non-branching rules add lines below; branching rules split into two branches (shown side by side).

\begin{multicols}{2}

\rulesubheader{Negation ($\neg$)}

\begin{tabular}{@{}c@{\qquad\qquad}c@{}}
$\T\, \neg A$ & $\F\, \neg A$ \\[2pt]
$\Downarrow$ & $\Downarrow$ \\[2pt]
$\F\, A$ & $\T\, A$
\end{tabular}

\smallskip
Both rules are non-branching.

\rulesubheader{Conjunction ($\wedge$)}

\begin{tabular}{@{}c@{\qquad\qquad}c@{}}
$\T\, A \wedge B$ & $\F\, A \wedge B$ \\[2pt]
$\Downarrow$ & $\swarrow\quad\searrow$ \\[2pt]
$\T\, A$ & $\F\, A \quad\mid\quad \F\, B$ \\
$\T\, B$ &
\end{tabular}

\smallskip
True: non-branching (both conjuncts true).\\
False: branching (at least one conjunct false).

\columnbreak

\rulesubheader{Disjunction ($\vee$)}

\begin{tabular}{@{}c@{\qquad\qquad}c@{}}
$\T\, A \vee B$ & $\F\, A \vee B$ \\[2pt]
$\swarrow\quad\searrow$ & $\Downarrow$ \\[2pt]
$\T\, A \quad\mid\quad \T\, B$ & $\F\, A$ \\
 & $\F\, B$
\end{tabular}

\smallskip
True: branching (at least one disjunct true).\\
False: non-branching (both disjuncts false).

\rulesubheader{Conditional ($\rightarrow$)}

\begin{tabular}{@{}c@{\qquad\qquad}c@{}}
$\T\, A \rightarrow B$ & $\F\, A \rightarrow B$ \\[2pt]
$\swarrow\quad\searrow$ & $\Downarrow$ \\[2pt]
$\F\, A \quad\mid\quad \T\, B$ & $\T\, A$ \\
 & $\F\, B$
\end{tabular}

\smallskip
True: branching (antecedent false or consequent true).\\
False: non-branching (antecedent true, consequent false).

\end{multicols}

\medskip
\textbf{Closure:} A branch closes when it contains both $\T\, A$ and $\F\, A$ for some sentence $A$. A tableau closes when \emph{every} branch closes.

\textbf{Strategy:} Apply non-branching rules before branching rules.

\bigskip
\ruleheader{Part 2: Modal Rules in K}

In \textbf{K}, worlds are numbered with dot notation: $x$, $x.y$, $x.y.z$, etc. World $x$ can access world $x.y$ (and only those forced by the dot structure).

\begin{multicols}{2}

\rulesubheader{True $\Box$ (open-ended, reusable)}

If $\T\, x{:}\; \Box A$ appears on the branch, then for \emph{every} world $x.y$ on the branch, add:
$$\T\, x.y{:}\; A$$
Apply this each time a new $x.y$ appears.

\rulesubheader{False $\Diamond$ (open-ended, reusable)}

If $\F\, x{:}\; \Diamond A$ appears on the branch, then for \emph{every} world $x.y$ on the branch, add:
$$\F\, x.y{:}\; A$$
Apply this each time a new $x.y$ appears.

\columnbreak

\rulesubheader{True $\Diamond$ (one-shot)}

If $\T\, x{:}\; \Diamond A$ appears on the branch, introduce a \emph{new} world $x.y$ (not yet on the branch) and add:
$$\T\, x.y{:}\; A$$

\rulesubheader{False $\Box$ (one-shot)}

If $\F\, x{:}\; \Box A$ appears on the branch, introduce a \emph{new} world $x.y$ (not yet on the branch) and add:
$$\F\, x.y{:}\; A$$

\end{multicols}

\smallskip
\textbf{Key distinction:} True $\Box$ / False $\Diamond$ are \emph{reusable}---they apply to every accessible world that appears. True $\Diamond$ / False $\Box$ are \emph{one-shot}---they introduce a single new world.

\newpage

\begin{center}
{\LARGE\bfseries Tableau Rules Cheat Sheet (cont.)}\\[4pt]
{\large Additional Rules for T, 4, and B}
\end{center}

\ruleheader{Part 3: Rules Beyond K}

These rules are \textbf{added to} the K rules. They only affect the open-ended rules (True $\Box$ and False $\Diamond$). The one-shot rules (True $\Diamond$ and False $\Box$) stay the same in all logics.

\bigskip

\rulesubheader{T Rules (Reflexivity: every world accesses itself)}

\begin{multicols}{2}

\textbf{T-Box Rule:} If $\T\, x{:}\; \Box A$ is on the branch, add:
$$\T\, x{:}\; A$$

(Since $x$ accesses itself, $A$ must hold at $x$ too.)

\columnbreak

\textbf{T-Diamond Rule:} If $\F\, x{:}\; \Diamond A$ is on the branch, add:
$$\F\, x{:}\; A$$

(Since $x$ accesses itself, $A$ must fail at $x$ too.)

\end{multicols}

\bigskip

\rulesubheader{4 Rules (Transitivity: if $x$ accesses $x.y$, then $x$ accesses everything $x.y$ accesses)}

\begin{multicols}{2}

\textbf{4-Box Rule:} If $\T\, x{:}\; \Box A$ is on the branch and $x.y$ is on the branch, add:
$$\T\, x.y{:}\; \Box A$$

(The whole $\Box A$, not just $A$, carries forward.)

\columnbreak

\textbf{4-Diamond Rule:} If $\F\, x{:}\; \Diamond A$ is on the branch and $x.y$ is on the branch, add:
$$\F\, x.y{:}\; \Diamond A$$

(The whole $\Diamond A$ false, not just $A$ false, carries forward.)

\end{multicols}

\bigskip

\rulesubheader{B Rules (Symmetry: if $x$ accesses $x.y$, then $x.y$ also accesses $x$)}

\begin{multicols}{2}

\textbf{B-Box Rule:} If $\T\, x.y{:}\; \Box A$ is on the branch, add:
$$\T\, x{:}\; A$$

(Since $x.y$ can look back to $x$, $A$ must hold at $x$.)

\columnbreak

\textbf{B-Diamond Rule:} If $\F\, x.y{:}\; \Diamond A$ is on the branch, add:
$$\F\, x{:}\; A$$

(Since $x.y$ can look back to $x$, $A$ must fail at $x$.)

\end{multicols}

\bigskip
\ruleheader{Part 4: Quick Reference --- Combining Logics}

\begin{center}
\begin{tabular}{lll}
\toprule
\textbf{Logic} & \textbf{Rules Used} & \textbf{Frame Condition} \\
\midrule
K & K rules only & No restrictions on $R$ \\
KT & K + T rules & Reflexive ($xRx$) \\
K4 & K + 4 rules & Transitive (if $xRy$ and $yRz$ then $xRz$) \\
KB & K + B rules & Symmetric (if $xRy$ then $yRx$) \\
S4 (= KT4) & K + T + 4 rules & Reflexive + Transitive \\
KTB & K + T + B rules & Reflexive + Symmetric \\
S5 (= KT4B) & K + T + 4 + B rules & Equivalence relation (= universal) \\
\bottomrule
\end{tabular}
\end{center}

\bigskip

\textbf{Characteristic axioms:}
\begin{itemize}[nosep]
\item \textbf{T:} $\Box A \rightarrow A$ \quad (what's necessary is actual)
\item \textbf{B:} $A \rightarrow \Box\Diamond A$ \quad (what's actual is necessarily possible)
\item \textbf{4:} $\Box A \rightarrow \Box\Box A$ \quad (what's necessary is necessarily necessary)
\end{itemize}

\bigskip

\textbf{Reminder for S5 (simplified tableaux):} In S5, every world accesses every world. So:
\begin{itemize}[nosep]
\item $\T\, \Box A$ at any world $\Rightarrow$ $\T\, A$ at \emph{every} world.
\item $\F\, \Diamond A$ at any world $\Rightarrow$ $\F\, A$ at \emph{every} world.
\item $\T\, \Diamond A$ at any world $\Rightarrow$ introduce one new world with $\T\, A$ (can't reuse a world where $A$ is already false).
\item $\F\, \Box A$ at any world $\Rightarrow$ introduce one new world with $\F\, A$ (can't reuse a world where $A$ is already true).
\end{itemize}

\end{document}
